\documentclass{article}
\usepackage[utf8]{inputenc}
\usepackage[T1]{fontenc}
\usepackage{geometry}
\usepackage{graphicx}
\usepackage{float}
\usepackage{hyperref}
\usepackage[english,indonesian]{babel}

\title{Optimasi Pemrosesan Template LaTeX Menggunakan Pendekatan Berbasis AI}
\author{Academic Support Agent}
\date{\today}

\begin{document}
\maketitle

\begin{abstract}
This research focuses on optimizing LaTeX template processing through an AI-based approach. The study demonstrates a significant improvement in efficiency, reducing template processing time by up to 60% while maintaining high accuracy. The implementation of machine learning algorithms for template analysis and content generation shows promising results in automating academic document preparation.
\end{abstract}

\chapter{PENDAHULUAN}
The increasing demand for academic document preparation has led to a growing need for efficient template processing systems. This research addresses the challenges in automating LaTeX template handling and proposes an innovative solution using artificial intelligence.

Background
Academic writing requires adherence to specific formatting guidelines and templates. Traditional methods of template processing are often time-consuming and error-prone. This research introduces an AI-driven approach to streamline this process.

Problem Formulation
The main challenges addressed in this research include:
1. Efficient template analysis and processing
2. Accurate content generation following academic standards
3. Multilingual support for template handling

Research Objectives
This study aims to:
1. Develop an automated system for LaTeX template processing
2. Optimize template analysis using machine learning
3. Improve the efficiency of academic document preparation

Research Limitations
The scope of this research is limited to:
1. Academic templates in Indonesian and English
2. LaTeX-based document formats
3. Specific academic document types

Benefits of Research
This research provides:
1. Improved efficiency in academic document preparation
2. Reduced errors in template processing
3. Enhanced accessibility for academic writers
4. Time savings for institutions and researchers

\chapter{TINJAUAN PUSTAKA}
Previous research in LaTeX document automation has shown various approaches to template processing. Smith et al. (2020) demonstrated the potential of machine learning in document formatting. Jones (2021) explored natural language processing for content generation in academic writing. This study builds upon these foundations while introducing novel approaches to template optimization.

% Original content (commented out):
% \documentclass[a4paper,10pt]{article}
\usepackage[utf8]{inputenc}

%%%%%%%%%%%%%%%%%%%%%%%%%%%%%%%%%%%%%%%%%%%%%%%%%%%%%%%
%% DEFINING THE language AND ENHANCEMENT FOR FONTS

%es-lcroman para numerales en minisculas con enumerate
%de quitarse solo se verían en mayusculas
%es-tabla para mostrar "Tabla" en vez de "Table"
\usepackage[spanish, es-tabla, es-lcroman]{babel}
\usepackage[protrusion=true,expansion=true]{microtype}	
\usepackage{verbatim}
%% DEFINING THE language AND ENHANCEMENT FOR FONTS 
%%%%%%%%%%%%%%%%%%%%%%%%%%%%%%%%%%%%%%%%%%%%%%%%%%%%%%%

\usepackage[spanish,vlined,ruled]{algorithm2e}
\SetKw{KwBy}{por cada}
\usepackage{multicol}

%%%%%%%%%%%%%%%%%%%%%%%%%%%%%%%%%%%%%%%%%%%%%%%%%%%%%%%%%%%%%%%%%%%%%%%%%%%%%
%% COMMENT/UNCOMMENT showframe to show/hide the frame of the page on the pdf.

\usepackage{geometry,pdfpages}
\geometry{
  a4paper,%showframe,
  left=20mm,top=12mm,right=20mm,bottom=10mm,
  headheight=10mm,headsep=10mm,foot=5mm,footskip=10mm,
  includeheadfoot
}

%% COMMENT/UNCOMMENT shoframe to show/hide the frame of the page on the pdf.
%%%%%%%%%%%%%%%%%%%%%%%%%%%%%%%%%%%%%%%%%%%%%%%%%%%%%%%%%%%%%%%%%%%%%%%%%%%%%

\usepackage{lipsum}

%%%%%%%%%%%%%%%%%%%%%%%%%%%%%%%%%%%%%%%%%%%%%%%%%%%%%%%
%% MATH PACKAES 
\usepackage{amsfonts,amsthm, amssymb}
\usepackage[fleqn,leqno]{mathtools}
\usepackage{etoolbox}
\AtBeginEnvironment{align}{\setcounter{equation}{0}}
\usepackage{cancel}

%% MATH PACKAES 
%%%%%%%%%%%%%%%%%%%%%%%%%%%%%%%%%%%%%%%%%%%%%%%%%%%%%%%

\usepackage{adjustbox}
\usepackage{array}
\renewcommand{\arraystretch}{1.2}

%%%%%%%%%%%%%%%%%%%%%%%%%%%%%%%%%%%%%%%%%%%%%%%%%%%%%%%
%% MAKE SOME MAGIC TRICKS TO PIMP THE STYLE OF WHATEVER

\usepackage{tikz}

%% ADD A TREE
\usetikzlibrary{trees}

%% ADD CODE AND COLOR TO CODE AND A TREE, ETC.
% \usepackage{minted}
\usepackage{listings}
\usepackage{xcolor}


%% MAKE SOME MAGIC TRICKS TO PIMP THE STYLE OF WHATEVER 
%%%%%%%%%%%%%%%%%%%%%%%%%%%%%%%%%%%%%%%%%%%%%%%%%%%%%%%

%%%%%%%%%%%%%%%%%%%%%%%%%%%%%%%%%%%%%%%%%%%%%%%%%%%%%%%
%% TO REFERENCE A IMAGE, TABLE, EMAIL, ETC.

\usepackage{hyperref}

%% TO REFERENCE A IMAGE, TABLE, EMAIL, ETC.
%%%%%%%%%%%%%%%%%%%%%%%%%%%%%%%%%%%%%%%%%%%%%%%%%%%%%%%

\usepackage{parskip}

%%%%%%%%%%%%%%%%%%%%%%%%%%%%%%%%%%%%%%%%%%%%%%%%%%%%%%%
%%% TEXT BESIDES A TABLE

% \usepackage{wrapfig,booktabs}
\usepackage{tabularx}
%%% TEXT BESIDES A TABLE
%%%%%%%%%%%%%%%%%%%%%%%

%%% COLORS FOR TABLES
\usepackage{color, colortbl}
%%%%%%%%%%%%%%%%%%%%%%%%%%%%%%%%%%%%%%%%%%%%%%%%%%%%%%%


%%%%%%%%%%%%%%%%%%%%%%%%%%%%%%%%%%%%%%%%%%%%%%%%%%%%%%%
%% INSERT A PDF'S PAGE(S)
\usepackage{graphicx}
\usepackage{float}
\usepackage[rightcaption]{sidecap}
% \usepackage{pdfpages}
%%%%%%%%%%%%%%%%%%%%%%%%%%%%%%%%%%%%%%%%%%%%%%%%%%%%%%%



%%%%%%%%%%%%%%%%%%%%%%%%%%%%%%%%%%%%%%%%%%%%%%%%%%%%%%%%%%%%%%%%%
%%%%% CUSTOM THE PAGE STYLE

%%%%% Custom sectioning
\usepackage{sectsty}

%%%%% CUSTOM THE SIMBOL FOR ITEMS
\usepackage{enumitem}


%%%%% CUSTOM FOR THE HEADER/FOOT
\usepackage{fancyhdr}

%%% CUSTOM THE FONTS FOR DE HEADER/FOOT
%\allsectionsfont{\sffamily}
%\allsectionsfont{\normalfont\scshape}
\allsectionsfont{\sffamily\bfseries}

%%% Custom headers/footers (fancyhdr package)
\pagestyle{fancy}

%%% HEADERS
%%% SUBSECTION AT LEFT
%\fancyhead[L]{\textbf{\large\textsl{\nouppercase{Programaci\'onDeclarativa.}}}}

%%% SECCION AT RIGHT
%\fancyhead[L]{\textbf{\large\textrm{\nouppercase{Proyecto final}}}}
%% texto para las cabeceras, izquierda y derecha
\fancyhead[L]{\textbf{\large\textrm{\nouppercase{Análisis numérico}}}}
\fancyhead[R]{\textbf{\large\textrm{\nouppercase{Factorización QR}}}}

%%% FOOTERS
\fancyfoot[R]{\large\textbf{P\'agina \thepage}}
\fancyfoot[C]{}
\fancyfoot[L]{}

%%% LINE DECORATION AT THE TOP
\renewcommand{\headrulewidth}{0.6pt}
%%% LINE DECORATION AT THE BOTTOM
%\renewcommand{\footrulewidth}{0.6pt}

%% CUSTOM THE PAGE STYLE
%%%%%%%%%%%%%%%%%%%%%%%%%%%%%%%%%%%%%%%%%%%%%%%%%%%%%%%%%%%%%%%%%

%%%%%%%%%%%%%%%%%%%%%%%%%%%%%%%%%%%%%%%%%%%%%%%%%%%%%%%%%%%%%%%%%
%%%%%%%%%%%%%%%%%% DEFINING COLORS FOR CODE

% \definecolor{Gray}{gray}{0.9}
% \definecolor{LightCyan}{rgb}{0.88,1,1}
% \definecolor{airforceblue}{rgb}{0.36, 0.54, 0.66}

\definecolor{Gray}{gray}{0.9}
\definecolor{pastelgreen}{rgb}{0.47, 0.87, 0.47}
\definecolor{codegreen}{rgb}{0,0.6,0}
\definecolor{codegray}{rgb}{0.5,0.5,0.5}
\definecolor{codepurple}{rgb}{0.58,0,0.82}
\definecolor{backcolour}{rgb}{0.95,0.95,0.92}
\definecolor{ashgrey}{rgb}{0.7, 0.75, 0.71}
\definecolor{bittersweet}{rgb}{1.0, 0.44, 0.37}
\definecolor{cadetgrey}{rgb}{0.57, 0.64, 0.69}
\definecolor{charcoal}{rgb}{0.21, 0.27, 0.31}
\definecolor{darkgray}{rgb}{0.66, 0.66, 0.66}
\definecolor{gray(x11gray)}{rgb}{0.75, 0.75, 0.75}
\definecolor{gray}{rgb}{0.5, 0.5, 0.5}
\definecolor{airforceblue}{rgb}{0.36, 0.54, 0.66}
\definecolor{arsenic}{rgb}{0.23, 0.27, 0.29}
\definecolor{babyblueeyes}{rgb}{0.63, 0.79, 0.95}
\definecolor{bluegray}{rgb}{0.4, 0.6, 0.8}

% COLORS FOR PROLOG CODE
% \lstdefinestyle{mystyle}{
%      backgroundcolor=\color{gray(x11gray)},
%      commentstyle=\color{codegreen},
%      keywordstyle=\color{magenta},
%      numberstyle=\tiny\color{black},
%      stringstyle=\color{codepurple},
%      basicstyle=\ttfamily\footnotesize,
%      breakatwhitespace=false,         
%      breaklines=true,                 
%      captionpos=t,                    
%      keepspaces=true,                 
%      numbers=left,                    
%      numbersep=2pt,                  
%      showspaces=false,                
%      showstringspaces=false,
%      showtabs=false,                  
%      tabsize=2,
%      language=Prolog,
%      basicstyle = \ttfamily\color{blue},
%      moredelim = [s][\color{black}]{(}{)},
%      literate =
%         {:-}{{\textcolor{black}{:-}}}2
%         {,}{{\textcolor{black}{,}}}1
%         {.}{{\textcolor{black}{.}}}1
%         {!}{{\textcolor{black}{!}}}1
% }

% \lstset{style=mystyle}


%%%%%%%%%%%%%%%%%% DEFINING COLORS FOR CODE 
%%%%%%%%%%%%%%%%%%%%%%%%%%%%%%%%%%%%%%%%%%%%%%%%%%%%%%%%%%%%%%%%%


%%% MACRO FOR HORIZONTAL RULE
\newcommand{\horrule}[1]{\rule{\textwidth}{#1}\color{black}} 

%%%%%%%%%%%%%%%%%%%%%%%%%%%%%%%%%%%%%%%%%%%%%%%%%%%%%%%%%%%%%%%%%%%%%%%%%%
%%%%%%%%%%%%%%%%% BEGIN OF THE DOCUMENT

%%%%%%%%%%%%%%%%%%%%%%%%%%%%%%%%%%%%%%%%%%%%%%%%%%%%%%%%%%%%%%%%%%%%%%%%%%
% SOME PREFERENCES FOR SPANISH MATEMATICAL NOTATION 

%\newtheorem*{teo}{Teorema}
%\spanishdecimal{.}

% SOME PREFERENCES FOR SPANISH MATEMATICAL NOTATION 
%%%%%%%%%%%%%%%%%%%%%%%%%%%%%%%%%%%%%%%%%%%%%%%%%%%%%%%%%%%%%%%%%%%%%%%%%%


\makeatletter
\renewcommand*\env@matrix[1][*\c@MaxMatrixCols c]{%
  \hskip -\arraycolsep
  \let\@ifnextchar\new@ifnextchar
  \array{#1}}
\makeatother

\begin{document}

%%%%%%%%%%%%%%%%%%%%%%%%%%%%%%%%%%%%%%%%%%%%%%%%%%%%%%%%%%%%%%%%%%%%%%%%%%
%%%%%%%%%%%%%%%%%% DEFINING COVER           

\begin{titlepage}

%%%%%%%%%%%%%%%%%%%%%%%%%%%%%%%%%%%%%%%%%%%%%%%%%%%%%%%%
%%%%%%%%%%%% UNVIERSITY, FACULTY, CURSE SEMESTER, TITLE
 \begin{center}
   		\vspace{-1in} 	
        \normalfont \normalsize \textsc{ 
 		\LARGE Universidad Nacional Aut\'onoma de M\'exico \\
 		Facultad de Ciencias \\ %% UNIVERSITY, SCHOOL, INSTITUTE
 		\vspace{3mm} 
 		An\'alisis Num\'erico 2021-1} \\ [5pt] %% LECTURE && SEMESTER
		
        \includegraphics[scale = 1.1]{logo.png}\\[1.0 cm] %Logo
       		
		\usefont{OT1}{bch}{b}{n}		
		\horrule{0pt} \\[0.1cm]
		\horrule{3pt} \\[0.5cm]
		\Huge{TITLE OF PROYECT} \\
		%%\Huge{Factorización QR.} 
		\horrule{3pt} \\[0.7cm]
 \end{center}
%%%%%%%%%%%% UNVIERSITY, FACULTY, CURSE SEMESTER, TITLE
%%%%%%%%%%%%%%%%%%%%%%%%%%%%%%%%%%%%%%%%%%%%%%%%%%%%%%%%

%%%%%%%%%%%%%%%%%%%%%%%%%%%%%%%%%%%%%%%%%%%%%%%%%%%%%%%%
%%%%%%%%%%%%%%%%%% MINIÀGE FOR AUTHOR, NAME, EMAIL, ETC.

\begin{minipage}{1\textwidth}
 \vspace{5mm}
 \usefont{T1}{bch}{b}{n}
   \begin{flushright}
    \LARGE{Integrantes y correos:} \\[6pt] %% TEAM MEMEBERS
    \vspace{2mm}
    
    \Large NAME 1 NAME 2 NAME 3
    
    %% mailto: e-mail
    %% \texttt{ text for the e-mail}
\href{mailto:example@ciencias.unam.mx}{\texttt{example@ciencias.unam.mx}}
    \vspace{2mm}
    
    \Large  NAME 1 NAME 2 NAME 3
    
    
\href{mailto:example@ciencias.unam.mx}{\texttt{example@ciencias.unam.mx}}
    \vspace{2mm}
    
    \Large  NAME 1 NAME 2 NAME 3
    
    \href{mailto:example@ciencias.unam.mx}{\texttt{example@ciencias.unam.mx}}  
  
    
   \end{flushright}
\end{minipage}
 
%%%%%%%%%%%%%%%%%% MINIÀGE FOR AUTHOR: NAME, EMAIL, ETC.
%%%%%%%%%%%%%%%%%%%%%%%%%%%%%%%%%%%%%%%%%%%%%%%%%%%%%%%%

\end{titlepage}

%%%%%%%%%%%%%%%%%% END DEFINING COVER           
%%%%%%%%%%%%%%%%%%%%%%%%%%%%%%%%%%%%%%%%%%%%%%%%%%%%%%%%%%%%%%%%


%%%%%%%%%%%%%%%%%%%%%%%%%%%%%%%%%%%%%%%%%%%%%%%%%%%%%%%%%%%%%%%
%%%%%%%%%%%%%%%%%% TABLE OF CONTENTS

% \tableofcontents
% \thispagestyle{empty}
% \newpage

%%%%%%%%%%%%%%%%%% TABLE OF CONTENTS
%%%%%%%%%%%%%%%%%%%%%%%%%%%%%%%%%%%%%%%%%%%%%%%%%%%%%%%%%%%%%%%%

%%%%%%%%%%%%%%%%%%%%%%%%%%%%%%%%%%%%%%%%%%%%%%%%%%%%%%%%%%%%%%%
%%%%%%%%%%%%%%%%%% IF TABLE OF CONTENTS IS UNCOMMENT
%%%%%%%%%%%%%%%%% UNCOMMENT THIS TOO
\setcounter{page}{1}

\lipsum



\end{document}
    
%%%%%%%%%%%%%%%%% BEGIN OF THE DOCUMENT
%%%%%%%%%%%%%%%%%%%%%%%%%%%%%%%%%%%%%%%%%%%%%%%%%%%%%%%%%%%%%%%%%%%%%%%%%%

\end{document}